% ══════════════════════════════════════════════════════════════════
% INTRODUCTION GÉNÉRALE — 6 sections (plan de travail)
% ══════════════════════════════════════════════════════════════════

\chapter*{Introduction Générale}
\addcontentsline{toc}{chapter}{Introduction Générale}
\markboth{Introduction Générale}{Introduction Générale}

% ── 1. Contexte général ───────────────────────────────────────────
\section*{1. Contexte général}
\addcontentsline{toc}{section}{1. Contexte général}

L'agriculture occupe une place centrale dans l'économie tchadienne :
elle contribue à près de 20\,\% du Produit Intérieur Brut et fait
vivre plus de 80\,\% de la population rurale
\cite{banquemondiale2023}. Les principales cultures du pays --- le
mil (\textit{Pennisetum glaucum}), le sorgho (\textit{Sorghum
	bicolor}), l'arachide (\textit{Arachis hypogaea}), le maïs
(\textit{Zea mays}) et le coton (\textit{Gossypium hirsutum}) ---
sont pratiquées majoritairement en régime pluvial, dans des
conditions climatiques difficiles propres aux zones sahélienne et
soudanienne : forte variabilité des précipitations, températures
élevées et épisodes de sécheresse récurrents.

\medskip

Parallèlement, l'essor de l'Intelligence Artificielle (IA) ouvre
des perspectives inédites pour l'agriculture de précision. Les
techniques d'apprentissage automatique permettent aujourd'hui
d'exploiter des volumes importants de données climatiques,
pédologiques et visuelles afin de produire des recommandations
agronomiques personnalisées et en temps réel
\cite{liakos2018machine}. Ces avancées représentent une opportunité
majeure pour les pays en développement, où l'encadrement agronomique
de proximité reste insuffisant.

% ── 2. Problématique ──────────────────────────────────────────────
\section*{2. Problématique}
\addcontentsline{toc}{section}{2. Problématique}

Malgré son poids économique, l'agriculture tchadienne demeure peu
productive. Les agriculteurs font face à plusieurs obstacles
majeurs : la rareté des agronomes qualifiés en zone rurale,
l'absence d'accès aux informations météorologiques en temps réel,
la difficulté à détecter précocement les maladies des cultures,
et une gestion souvent inadaptée des intrants --- eau d'irrigation
et engrais --- entraînant des pertes de rendement et des coûts
inutiles.

\medskip

Face à ce constat, la question centrale de ce travail est la
suivante : \textbf{comment concevoir un outil d'aide à la décision
	accessible, précis et adapté aux conditions locales, capable
	d'accompagner l'agriculteur tchadien dans le suivi quotidien de
	ses parcelles ?}

\medskip

Cette problématique se décline en quatre questions opérationnelles :

\begin{itemize}[leftmargin=1.2cm]
	\item Comment prédire le besoin en irrigation d'une parcelle à
	partir des conditions climatiques locales en temps réel ?
	\item Comment recommander l'engrais optimal selon les
	caractéristiques du sol (NPK, pH) et le climat ?
	\item Comment détecter automatiquement les maladies des cultures
	à partir d'une photographie de feuille ou d'un indice de
	végétation ?
	\item Comment assurer la traçabilité des analyses effectuées
	afin de permettre un suivi historique des parcelles ?
\end{itemize}

% ── 3. Objectifs du travail ───────────────────────────────────────
\section*{3. Objectifs du travail}
\addcontentsline{toc}{section}{3. Objectifs du travail}

L'objectif général de ce projet est de \textbf{concevoir et
	d'implémenter un cadre décisionnel basé sur l'Intelligence
	Artificielle pour l'optimisation du suivi parcellaire agricole au
	Tchad}, matérialisé par une application web dénommée
\textbf{Agro-IA Tchad}.

\medskip

Cet objectif général se décompose en cinq objectifs spécifiques :

\begin{enumerate}[leftmargin=1.2cm]
	\item Développer un \textbf{module de prédiction du besoin en
		irrigation} fondé sur l'algorithme Random Forest, entraîné
	sur des données climatiques réelles de la base NASA POWER
	couvrant cinq villes tchadiennes sur la période 2015--2023 ;
	\item Concevoir un \textbf{module de recommandation d'engrais}
	utilisant l'algorithme XGBoost à partir des paramètres NPK
	du sol et des conditions climatiques ;
	\item Implémenter un \textbf{module de détection des maladies}
	combinant un réseau de neurones convolutif (CNN MobileNetV2)
	appliqué aux photographies de feuilles, et l'analyse
	d'indices de végétation NDVI obtenus par satellite
	Sentinel-2 ou estimés par caméra smartphone ;
	\item Mettre en place un \textbf{module Historique} assurant la
	persistance de toutes les analyses dans une base de données
	SQLite, avec statistiques et export CSV ;
	\item \textbf{Déployer l'application} sur une plateforme cloud
	accessible gratuitement depuis tout navigateur web.
\end{enumerate}

% ── 4. Motivation et intérêt du sujet ─────────────────────────────
\section*{4. Motivation et intérêt du sujet}
\addcontentsline{toc}{section}{4. Motivation et intérêt du sujet}

Le choix de ce sujet est motivé par un triple intérêt.

\medskip

\textbf{Sur le plan scientifique}, ce travail explore l'application
concrète de plusieurs familles d'algorithmes d'apprentissage
automatique --- forêts aléatoires, gradient boosting et réseaux de
neurones convolutifs --- à des données agricoles réelles et
hétérogènes, dans un contexte de ressources limitées (peu de
données annotées localement).

\medskip

\textbf{Sur le plan socio-économique}, un outil d'aide à la décision
gratuit et accessible peut contribuer à améliorer les rendements
agricoles, à réduire le gaspillage d'eau et d'engrais, et à limiter
les pertes de récoltes dues aux maladies, avec un impact direct sur
la sécurité alimentaire des ménages ruraux tchadiens.

\medskip

\textbf{Sur le plan personnel}, ce projet est l'occasion de mettre
en \oe uvre l'ensemble des compétences acquises durant la formation
de Licence à l'ENASTIC --- programmation, bases de données, génie
logiciel et intelligence artificielle --- au service d'une
problématique nationale concrète.

% ── 5. Méthodologie adoptée ───────────────────────────────────────
\section*{5. Méthodologie adoptée}
\addcontentsline{toc}{section}{5. Méthodologie adoptée}

Pour atteindre les objectifs fixés, une démarche en quatre phases
a été suivie.

\medskip

La \textbf{première phase} a été consacrée à la collecte et au
prétraitement des données : extraction de 16~435 observations
climatiques depuis l'API NASA POWER, constitution d'un jeu de
données de fertilisation de 1~599 observations, et création d'un
dataset local de 287 images de feuilles de cultures tchadiennes
réparties en 10 classes.

\medskip

La \textbf{deuxième phase} a porté sur l'entraînement et
l'optimisation des modèles d'IA sur la plateforme Kaggle (GPU
Tesla T4) : Random Forest pour l'irrigation, XGBoost pour la
fertilisation et CNN MobileNetV2 par transfert d'apprentissage
pour la détection des maladies, avec conversion finale au format
TensorFlow Lite.

\medskip

La \textbf{troisième phase} a concerné le développement de
l'application web avec le framework Streamlit : intégration des
modèles, de l'API météorologique Open-Meteo et de la base de
données SQLite pour le module Historique.

\medskip

Enfin, la \textbf{quatrième phase} a été dédiée aux tests,
à l'évaluation des performances de chaque module et au déploiement
de l'application sur Streamlit Cloud.

% ── 6. Structure du mémoire ───────────────────────────────────────
\section*{6. Structure du mémoire}
\addcontentsline{toc}{section}{6. Structure du mémoire}

Outre la présente introduction et la conclusion générale, ce
mémoire s'articule autour de quatre chapitres.

\medskip

Le \textbf{Chapitre 1 --- Travaux existants} présente l'agriculture
de précision, la nature des données agricoles, les techniques d'IA
appliquées à l'agriculture, ainsi qu'une analyse des solutions
existantes permettant de positionner notre contribution.

\medskip

Le \textbf{Chapitre 2 --- Analyse, Conception et Modélisation UML}
expose l'analyse du problème, la spécification des besoins, la
modélisation UML du système, l'architecture du cadre décisionnel
et le modèle de la base de données de l'historique.

\medskip

Le \textbf{Chapitre 3 --- Implémentation et Tests} détaille
l'environnement de développement, le prétraitement des données,
la réalisation des trois modules d'IA, du module Historique et de
l'application web, puis définit le protocole de tests et les
métriques d'évaluation.

\medskip

Le \textbf{Chapitre 4 --- Résultats et Discussion} présente les
résultats obtenus pour chaque module, la validation du module
Historique, une discussion générale et le déploiement de
l'application.

\medskip

Le mémoire s'achève par une \textbf{conclusion générale} qui
dresse la synthèse des travaux, en expose les limites et propose
des perspectives d'amélioration.